
% ================================================================
% Bridge Lemma (WKB-Überdeckung / Geometrische Abdeckung)
%
% Dieses File kann direkt in paper5.tex als neuer Abschnitt
% eingefügt werden. Alle Label-Präfixe beginnen mit "bridge:".
%
% Voraussetzungen (bereits in paper5.tex vorhanden):
%   - Lemma lem:eigenvalue-decay  (Eigenvalue concentration)
%   - Notation: psi_n, lambda_n, T, omega = c/T, N_Sh = 2c/pi
%   - Zitierbar: [OsipovRokhlinXiao], insb. §4.4 (WKB forbidden zone)
%
% Position im Paper: nach Section 3 (Off-diagonal decay),
%                    vor Section 4 (Schur argument).
% ================================================================

% ================================================================
\section{The Bridge Lemma: WKB Cover and Exponential Out-of-Band
Decay}
\label{sec:bridge}
% ================================================================

This section proves the key energy bound that closes the connection
between the WKB tunnel structure and the out-of-band spectral energy
of PSWF products. It establishes Assumption~2.4 of \cite{PaperIII}
unconditionally in the regime $\gamma < 1/2$, thereby completing the
bulk stability program for that parameter range.

% ---------------------------------------------------------------
\subsection{Setup and statement}
\label{subsec:bridge-setup}
% ---------------------------------------------------------------

Fix $0 < \gamma' < \gamma < 1/2$.
For $m, n \le \gamma' N_{\mathrm{Sh}}$, define the product function
\begin{equation}
\label{eq:bridge:fmn}
  f_{mn}(x) := \psi_m(x)\,\psi_n(x),
  \qquad x \in [-T, T].
\end{equation}
Since each $\psi_n \in L^2(\mathbb{R})$ is $\omega$-bandlimited
($\omega := c/T$), the product $f_{mn}$ lies in
$L^2(\mathbb{R})$ with Fourier support in $[-2\omega, 2\omega]$.
The \emph{out-of-band energy} of $f_{mn}$ is
\begin{equation}
\label{eq:bridge:Eout}
  \mathcal{E}_{\mathrm{out}}(f_{mn})
  := \int_{|\xi| > \omega} |\widehat{f}_{mn}(\xi)|^2\,d\xi.
\end{equation}

\begin{theorem}[Bridge Lemma: exponential out-of-band decay]
\label{thm:bridge}
Fix $0 < \gamma' < \gamma < 1/2$.
There exist constants $C_{\gamma,\gamma'} > 0$ and
$K_{\gamma,\gamma'} > 0$, depending only on $\gamma$ and $\gamma'$,
such that for all $m, n \le \gamma'\,N_{\mathrm{Sh}}$ and all
$c \ge 1$:
\begin{equation}
\label{eq:bridge:main}
  \mathcal{E}_{\mathrm{out}}(f_{mn})
  \;\le\;
  C_{\gamma,\gamma'}\,e^{-K_{\gamma,\gamma'}\,c}.
\end{equation}
In particular, Assumption~\textup{2.4} of \cite{PaperIII}
\textup{(ass:bulkconv)} holds unconditionally for all
$m, n \le \gamma'\,N_{\mathrm{Sh}}$ with $\gamma' < 1/2$.
\end{theorem}

% ---------------------------------------------------------------
\subsection{The WKB forbidden-zone pointwise bound}
\label{subsec:bridge-wkb}
% ---------------------------------------------------------------

The geometric heart of the argument rests on a pointwise bound for
PSWFs in their classical forbidden zone. We state this as a
stand-alone lemma.

\begin{lemma}[WKB forbidden-zone bound]
\label{lem:bridge:wkb}
Fix $0 < \gamma' < \gamma < 1$.
There exist constants $C_{\gamma,\gamma'} > 0$ and
$K_{\gamma,\gamma'} > 0$ such that for all
$n \le \gamma'\,N_{\mathrm{Sh}}$, all $x \in [\gamma T, T]$,
and all $c \ge 1$:
\begin{equation}
\label{eq:bridge:wkb}
  |\psi_n(x)|
  \;\le\;
  C_{\gamma,\gamma'}\,e^{-K_{\gamma,\gamma'}\,c}.
\end{equation}
\end{lemma}

\begin{proof}
\textbf{Step 1 (Turning point location).}
For $n \le \gamma' N_{\mathrm{Sh}} = 2\gamma' c/\pi$,
the WKB turning point $x_+(n)$ satisfies
$x_+(n) \le \gamma' T$ (see
\cite[§4.4, eq.~(4.85)]{OsipovRokhlinXiao}).
In particular, the entire interval $[\gamma T, T]$ lies in the
classically forbidden zone for $\psi_n$.

\textbf{Step 2 (WKB action lower bound).}
The WKB tunneling integral from $\gamma' T$ to $\gamma T$ is
\begin{equation}
\label{eq:bridge:I}
  I_{\gamma,\gamma'}
  := \int_{\gamma' T}^{\gamma T}
     \sqrt{\frac{c^2 t^2 - \chi_n(c)\,T^2}{T^2(T^2 - t^2)}}\,
     \frac{dt}{T}.
\end{equation}
Since $n \le \gamma' N_{\mathrm{Sh}}$ implies
$\chi_n(c) \le (c\,\gamma')^2 + C\,c$ for a uniform constant $C$
(cf.\ \cite[Theorem~2.7]{OsipovRokhlinXiao}), and since
$c^2 t^2 \ge c^2 (\gamma' T)^2$ for $t \ge \gamma' T$,
the integrand is bounded below uniformly:
\begin{equation}
\label{eq:bridge:integrand-lb}
  \frac{c^2 t^2 - \chi_n(c)\,T^2}{T^2(T^2 - t^2)}
  \;\ge\;
  \frac{c^2(\gamma^2 - (\gamma')^2)T^2/2}{T^2(1-\gamma'^2)T^2}
  \;=:\;
  \frac{c^2\,\mu_{\gamma,\gamma'}}{T^2}
\end{equation}
for all $t \in [\gamma'T, \gamma T]$ and all $c$ large enough,
where $\mu_{\gamma,\gamma'} := (\gamma^2-(\gamma')^2)/(2(1-(\gamma')^2)) > 0$.
Therefore
\begin{equation}
\label{eq:bridge:I-lower}
  I_{\gamma,\gamma'}
  \;\ge\;
  c\,\sqrt{\mu_{\gamma,\gamma'}}\,(\gamma - \gamma')
  \;=:\;
  c\,J_{\gamma,\gamma'}\;>\; 0.
\end{equation}

\textbf{Step 3 (WKB pointwise bound).}
By the uniform WKB asymptotic for PSWF eigenfunctions in the
forbidden zone
\cite[Theorem~4.6, eq.~(4.131)]{OsipovRokhlinXiao},
for $x \in [\gamma T, T - \varepsilon]$:
\[
  |\psi_n(x)| \;\le\;
  C\,(1 - x^2/T^2)^{-1/4}\,
  \exp\!\Bigl(-\int_{x_+(n)}^{x}
    \sqrt{V_c(t)}\,dt\Bigr),
\]
where $V_c(t) := (c^2t^2 - \chi_n T^2)/(T^2(T^2 - t^2))$.
The exponential factor satisfies
\[
  \exp\!\Bigl(-\int_{x_+(n)}^{x}\sqrt{V_c(t)}\,dt\Bigr)
  \;\le\;
  \exp\!\Bigl(-\int_{x_+(n)}^{\gamma T}\sqrt{V_c(t)}\,dt\Bigr)
  \;\le\;
  e^{-I_{\gamma,\gamma'}}
  \;\le\;
  e^{-c\,J_{\gamma,\gamma'}}.
\]
The prefactor $(1 - x^2/T^2)^{-1/4}$ is uniformly bounded on
$[\gamma T, T - \varepsilon]$.
Near $x = T$: the SL equation $(p(x)\psi_n')'
= (c^2x^2/T^2 - \chi_n)\psi_n$ with $p(T) = 0$ admits the
Frobenius analysis $\psi_n(x) = \psi_n(T) + O((T-x)^1)$,
with $|\psi_n(T)|$ bounded via the L$^2$-concentration bound
$\|\psi_n\|_{L^2[\gamma T, T]}^2 \le 1 - \lambda_n(c) \le e^{-c_1 c}$
(Lemma~\ref{lem:eigenvalue-decay}(i)) together with the Sobolev
interpolation
$|\psi_n(T)|^2 \le C\,\|\psi_n\|_{L^2}\,\|\psi_n'\|_{L^2}$
and the SL eigenvalue bound
$\|\psi_n'\|_{L^2}^2 \le \chi_n \le Cc^2$.
This gives $|\psi_n(T)| \le C\,c\,e^{-c_1 c/2}$.

Combining and setting $K_{\gamma,\gamma'} :=
\min(J_{\gamma,\gamma'},\, c_1/2)/2$:
\[
  |\psi_n(x)| \;\le\; C_{\gamma,\gamma'}\,e^{-K_{\gamma,\gamma'}\,c}
  \qquad\text{for all } x \in [\gamma T, T]. \qedhere
\]
\end{proof}

\begin{remark}[The endpoint $x = T$]
\label{rem:bridge:endpoint}
The SL operator is singular at $x = T$ because $p(T) = 0$.
However, the PSWFs are smooth up to the boundary, and the
pointwise value $\psi_n(T)$ is exponentially small: it satisfies
$|\psi_n(T)| \le C c\,e^{-c_1 c/2}$.
The polynomial prefactor $c$ is absorbed into the exponential
for the purposes of the Bridge Lemma.
\end{remark}

% ---------------------------------------------------------------
\subsection{The geometric cover argument}
\label{subsec:bridge-geo}
% ---------------------------------------------------------------

The core structural insight connecting WKB decay to
out-of-band Fourier energy is the following geometric lemma.

\begin{lemma}[Geometric forbidden-zone cover]
\label{lem:bridge:cover}
Fix $0 < \gamma < 1/2$.
Let $m, n \le \gamma N_{\mathrm{Sh}}$ and let
$\delta \in (0, \omega)$.
For every $x$ in the integration range $[\delta T/\omega, T]$,
at least one of the following holds:
\begin{enumerate}
  \item[\textup{(I)}] $x \ge \gamma T$, so
    $\psi_n(x)$ lies in the forbidden zone of $\psi_n$.
  \item[\textup{(II)}] $x < \gamma T$, so
    $T + \delta T/\omega - x \ge (1-\gamma)T > \gamma T \ge x_+(m)$,
    hence $\psi_m(T + \delta T/\omega - x)$ lies in the
    forbidden zone of $\psi_m$.
\end{enumerate}
In other words, for $\gamma < 1/2$, the two forbidden intervals
$(\gamma T, T)$ \textup{(}for $\psi_n$\textup{)} and the
pre-image interval $\{x : T + \delta T/\omega - x > \gamma T\} =
(-\infty, (1-\gamma)T + \delta T/\omega)$ \textup{(}for
$\psi_m$\textup{)} together cover all of $[\delta T/\omega, T]$.
\end{lemma}

\begin{proof}
For any $x \in [\delta T/\omega, T]$:
\begin{itemize}
  \item If $x \ge \gamma T$: Case~(I) applies.
  \item If $x < \gamma T$: then
    $T + \delta T/\omega - x > T + 0 - \gamma T = (1-\gamma)T$.
    Since $m \le \gamma N_{\mathrm{Sh}}$, the turning point satisfies
    $x_+(m) \le \gamma T$.
    Because $\gamma < 1/2$, we have $(1-\gamma) > \gamma$, so
    \[
      T + \delta T/\omega - x > (1-\gamma)T > \gamma T \ge x_+(m).
    \]
    Hence Case~(II) applies. \qedhere
\end{itemize}
\end{proof}

\begin{remark}[Why $\gamma < 1/2$ is sharp]
\label{rem:bridge:halfsharp}
The condition $\gamma < 1/2$ is the exact threshold for the
geometric cover to work.
For $\gamma = 1/2$: the two forbidden zones meet at the single
point $x = \gamma T = (1-\gamma)T = T/2$; the cover still holds
but with a degenerate boundary.
For $\gamma > 1/2$: there exists an interval
$x \in (\gamma T, (1-\gamma)T)$ in which
\emph{neither} factor is in its forbidden zone, and the
cover argument fails.
In that regime, a separate phase-cancellation argument
(as in Paper~I) is required.
\end{remark}

% ---------------------------------------------------------------
\subsection{Proof of the Bridge Lemma}
\label{subsec:bridge-proof}
% ---------------------------------------------------------------

\begin{proof}[Proof of Theorem~\ref{thm:bridge}]
By symmetry ($\mathcal{E}_{\mathrm{out}}(f_{mn}) =
\mathcal{E}_{\mathrm{out}}(f_{nm})$) we may assume $m, n \ge 0$.

\textbf{Step 1 (Convolution representation).}
Since $\psi_m, \psi_n$ are $\omega$-bandlimited and supported on
$[-T, T]$, the Fourier transform of $f_{mn}$ satisfies
\[
  \widehat{f}_{mn}(\xi)
  = (\widehat{\psi}_m * \widehat{\psi}_n)(\xi)
  = \int_{-\omega}^{\omega}
    \widehat{\psi}_m(\xi - \eta)\,\widehat{\psi}_n(\eta)\,d\eta.
\]
For $\xi = \omega + \delta$ with $\delta \in (0, \omega)$,
the integrand is supported on $\eta \in (\delta, \omega)$
(since $\widehat{\psi}_n$ is supported on $[-\omega, \omega]$ and
$\xi - \eta \in (\delta, \omega)$ requires $\eta \in (\delta, \omega)$).

Using the PSWF self-duality relation
$\widehat{\psi}_n(\eta) = (-i)^n \sqrt{\lambda_n}\,\psi_n(\eta T/\omega)$
for $|\eta| \le \omega$
(cf.\ \cite[§2.3, eq.~(2.37)]{OsipovRokhlinXiao}),
and the substitution $x = \eta T/\omega$:
\begin{equation}
\label{eq:bridge:conv-x}
  \widehat{f}_{mn}(\omega + \delta)
  = (-i)^{m+n}\sqrt{\lambda_m \lambda_n}\,
    \frac{\omega}{T}
    \int_{\delta T/\omega}^{T}
    \psi_m\!\bigl(T + \delta T/\omega - x\bigr)\,
    \psi_n(x)\,dx.
\end{equation}

\textbf{Step 2 (Pointwise bound via geometric cover).}
By Lemma~\ref{lem:bridge:cover} (with $\gamma'$ in
Lemma~\ref{lem:bridge:wkb} chosen as $\gamma' \to \gamma$),
for every $x \in [\delta T/\omega, T]$,
at least one of $|\psi_n(x)|$ or
$|\psi_m(T + \delta T/\omega - x)|$ is bounded by
$C_{\gamma}\,e^{-K_\gamma c}$ (Lemma~\ref{lem:bridge:wkb}),
while the other factor is bounded by $C$ (uniform $L^\infty$
bound for PSWFs; see \cite[Corollary~4.3]{OsipovRokhlinXiao}).
Therefore the integrand satisfies:
\[
  |\psi_m(T + \delta T/\omega - x)|\,|\psi_n(x)|
  \;\le\;
  C_\gamma\,e^{-K_\gamma c}
  \qquad\forall\, x \in [\delta T/\omega, T].
\]

\textbf{Step 3 (Integration).}
Integrating over $x \in [\delta T/\omega, T]$ (length $\le T$)
and using $|\sqrt{\lambda_m \lambda_n}| \le 1$:
\begin{equation}
\label{eq:bridge:pointwise-xi}
  |\widehat{f}_{mn}(\omega + \delta)|
  \;\le\;
  \frac{\omega}{T} \cdot T \cdot C_\gamma\,e^{-K_\gamma c}
  = C_\gamma\,\omega\,e^{-K_\gamma c}
  = C_\gamma\,\frac{c}{T}\,e^{-K_\gamma c}.
\end{equation}
By symmetry, the same bound holds for $\xi = -\omega - \delta$.

\textbf{Step 4 ($\mathcal{E}_{\mathrm{out}}$ bound).}
Integrating the squared bound over $\delta \in (0, \omega)$:
\begin{align}
\label{eq:bridge:Eout-final}
  \mathcal{E}_{\mathrm{out}}(f_{mn})
  &= 2\int_0^\omega |\widehat{f}_{mn}(\omega+\delta)|^2\,d\delta
  \;\le\;
  2\int_0^\omega C_\gamma^2\,\frac{c^2}{T^2}\,e^{-2K_\gamma c}\,d\delta
  \notag\\
  &= 2\,C_\gamma^2\,\frac{c^2}{T^2}\,\omega\,e^{-2K_\gamma c}
  = 2\,C_\gamma^2\,\frac{c^3}{T^3}\,e^{-2K_\gamma c}.
\end{align}
Since $c^3 e^{-2K_\gamma c} \le e^{-(2K_\gamma - \varepsilon)c}$
for any $\varepsilon > 0$ and $c$ large enough, setting
$K_{\gamma,\gamma'} := K_\gamma - \varepsilon/2$ yields
\begin{equation}
  \mathcal{E}_{\mathrm{out}}(f_{mn})
  \;\le\;
  C_{\gamma,\gamma'}\,e^{-K_{\gamma,\gamma'}\,c},
\end{equation}
which is \eqref{eq:bridge:main}. \qed
\end{proof}

% ---------------------------------------------------------------
\subsection{Consequences for the bulk stability program}
\label{subsec:bridge-consequences}
% ---------------------------------------------------------------

\begin{corollary}[ass:bulkconv is proved for $\gamma < 1/2$]
\label{cor:bridge:bulkconv}
For all $0 < \gamma < 1/2$, Assumption~\textup{ass:bulkconv}
of \cite{PaperIII} holds unconditionally:
\[
  \mathcal{E}_{\mathrm{out}}(f_{mn}) \le C_\gamma\,e^{-\alpha_\gamma c}
  \qquad \text{for all } m,n \le \gamma N_{\mathrm{Sh}},
\]
with $\alpha_\gamma = K_{\gamma,\gamma'} > 0$ for any
$\gamma' \in (\gamma/2,\gamma)$.
By \cite[Theorem~3.2 (thm:bulk)]{PaperIII}, the uniform
spectral tail bound
\[
  \|(I - P_N)f_{mn}\|_{L^2} \le C_\gamma\,e^{-\alpha_\gamma' N}
\]
follows unconditionally for all $m,n,N \le \gamma N_{\mathrm{Sh}}$
and $\gamma < 1/2$.
\end{corollary}

\begin{remark}[Separation of mechanisms]
\label{rem:bridge:separation}
The Bridge Lemma uses only two ingredients:
\begin{enumerate}
\item[(i)] WKB tunneling: $|\psi_n(x)| \le C e^{-Kc}$ in the
  forbidden zone.
\item[(ii)] Geometric cover: the two forbidden zones together
  cover the full convolution interval for $\gamma < 1/2$.
\end{enumerate}
In particular, \emph{no phase cancellation} (integration by
parts, stationary phase, or off-diagonal index decay in $|m-n|$)
is used.
The off-diagonal structure $|E_{mn}| \le C(1+|m-n|)^{-p}$
required for DSTP in \cite{PaperI} is a strictly additional
statement, established there by a separate phase-gap mechanism.
The Bridge Lemma proves only the energy bound; DSTP requires the
index-dependent oscillation control of \cite{PaperI}.
\end{remark}

\begin{remark}[Extension to $\gamma \ge 1/2$]
\label{rem:bridge:extension}
For $\gamma \ge 1/2$, the geometric cover fails: there exist
points $x \in (\gamma T, (1-\gamma)T)$ at which both
$\psi_n(x)$ and $\psi_m(T-x)$ are in the oscillatory regime.
In that zone, one expects cancellation from the WKB phase
$\Phi_{mn}(x) = \phi_m(x_0 - x) \pm \phi_n(x)$, which oscillates
at a rate governed by the index difference $|m-n|$ and the
spectral parameters $\chi_m, \chi_n$.
This leads to the stationary-phase / integration-by-parts
approach, which we do not pursue here.
Whether the exponential bound
$\mathcal{E}_{\mathrm{out}}(f_{mn}) \le C e^{-\alpha c}$
extends to $\gamma \in [1/2, 1)$ remains open.
\end{remark}
